\documentclass{beamer}

\usetheme{Madrid}

\usepackage{amsmath,amssymb,amsfonts}
\usepackage{bm}
\usepackage{graphicx}

\title{Powell's Iterative Method}
\subtitle{Derivative-Free Multidimensional Optimization}
\author{Morten Hjorth-Jensen}
\date{Spring 2026}

\begin{document}

\frame{\titlepage}

%------------------------------------------------
\section{Introduction}

\begin{frame}{Optimization in Computational Physics}

Many problems reduce to minimizing a function

\[
f(\mathbf{x}), \quad \mathbf{x} \in \mathbb{R}^n
\]

Examples:

\begin{itemize}
\item Energy minimization in quantum many-body problems
\item Fitting model parameters
\item Variational wavefunction optimization
\item Machine learning loss functions
\end{itemize}

Two major classes of methods:

\begin{itemize}
\item Gradient-based methods
\item Derivative-free methods
\end{itemize}

Powell's method belongs to the second class.
\end{frame}

%------------------------------------------------
\begin{frame}{Why Derivative-Free Methods?}

Gradients may be unavailable or expensive:

\begin{itemize}
\item Monte Carlo noise
\item Complex simulations
\item Black-box models
\end{itemize}

Derivative-free methods include:

\begin{itemize}
\item Nelder-Mead simplex
\item Powell's method
\item Pattern search
\end{itemize}

Powell's method is particularly efficient for smooth functions.
\end{frame}

%------------------------------------------------
\section{Basic Idea}

\begin{frame}{Concept of Direction Set Methods}

Goal: minimize

\[
f(\mathbf{x})
\]

Iteratively perform

\textbf{line minimizations}

along chosen directions.

Starting point

\[
\mathbf{x}_0
\]

Move along direction

\[
\mathbf{d}
\]

by minimizing

\[
\phi(\alpha) = f(\mathbf{x}_0 + \alpha \mathbf{d})
\]

This reduces a multidimensional problem to a sequence of 1D problems.
\end{frame}

%------------------------------------------------
\begin{frame}{Line Minimization}

For a given direction $\mathbf{d}$:

\[
\phi(\alpha) = f(\mathbf{x} + \alpha \mathbf{d})
\]

Goal:

\[
\alpha^* = \arg\min_\alpha \phi(\alpha)
\]

Methods:

\begin{itemize}
\item Golden section search
\item Brent's method
\end{itemize}

Update

\[
\mathbf{x}_{new} = \mathbf{x} + \alpha^* \mathbf{d}
\]
\end{frame}

%------------------------------------------------
\section{Powell's Method}

\begin{frame}{Direction Set}

Powell's method maintains a set of directions

\[
\{\mathbf{d}_1,\mathbf{d}_2,\dots,\mathbf{d}_n\}
\]

Initially chosen as unit vectors:

\[
\mathbf{d}_1 = (1,0,0,\dots)
\]

\[
\mathbf{d}_2 = (0,1,0,\dots)
\]

etc.
\end{frame}

%------------------------------------------------
\begin{frame}{One Iteration}

Starting at point $\mathbf{x}_0$.

Perform line minimizations sequentially:

\[
\mathbf{x}_1 = \mathbf{x}_0 + \alpha_1 \mathbf{d}_1
\]

\[
\mathbf{x}_2 = \mathbf{x}_1 + \alpha_2 \mathbf{d}_2
\]

\[
\dots
\]

\[
\mathbf{x}_n = \mathbf{x}_{n-1} + \alpha_n \mathbf{d}_n
\]

Each step minimizes along a direction.
\end{frame}

%------------------------------------------------
\begin{frame}{New Direction}

After completing the cycle define

\[
\mathbf{d}_{new} =
\mathbf{x}_n - \mathbf{x}_0
\]

This direction approximates the direction of largest improvement.

Replace one of the existing directions with $\mathbf{d}_{new}$.
\end{frame}

%------------------------------------------------
\section{Derivation and Motivation}

\begin{frame}{Quadratic Functions}

Consider quadratic function

\[
f(\mathbf{x}) =
\frac12 \mathbf{x}^T A \mathbf{x}
+ \mathbf{b}^T \mathbf{x} + c
\]

For such functions

\begin{itemize}
\item conjugate directions allow exact minimization
\item minimum found in $n$ steps
\end{itemize}

Powell's method approximates conjugate directions without gradients.
\end{frame}

%------------------------------------------------
\begin{frame}{Conjugate Directions}

Two directions $\mathbf{d}_i$ and $\mathbf{d}_j$ are conjugate if

\[
\mathbf{d}_i^T A \mathbf{d}_j = 0
\]

For quadratic functions this ensures

successive minimizations do not undo previous ones.

Powell's algorithm attempts to generate such directions adaptively.
\end{frame}

%------------------------------------------------
\section{Full Algorithm}

\begin{frame}{Powell Iteration}

Given

\[
\mathbf{x}_0
\]

and directions

\[
\mathbf{d}_1,\dots,\mathbf{d}_n
\]

\begin{enumerate}
\item Perform line minimizations along each direction
\item Record largest decrease in function
\item Construct new direction
\[
\mathbf{d}_{new} = \mathbf{x}_n - \mathbf{x}_0
\]
\item Replace direction with largest decrease
\item Repeat
\end{enumerate}
\end{frame}

%------------------------------------------------
\begin{frame}{Stopping Criteria}

Terminate when

\begin{itemize}
\item function change small
\[
|f_{new}-f_{old}| < \epsilon
\]

\item step size small

\[
||\mathbf{x}_{new}-\mathbf{x}_{old}|| < \epsilon
\]

\item maximum iterations reached
\end{itemize}
\end{frame}

%------------------------------------------------
\section{Numerical Aspects}

\begin{frame}{Advantages}

\begin{itemize}
\item No derivatives required
\item Efficient for moderate dimensions
\item Works well for smooth functions
\end{itemize}
\end{frame}

%------------------------------------------------
\begin{frame}{Limitations}

\begin{itemize}
\item Not robust for noisy functions
\item May fail for non-smooth problems
\item Performance deteriorates in high dimensions
\end{itemize}
\end{frame}

%------------------------------------------------
\section{Pseudocode}

\begin{frame}[fragile]{Powell Algorithm}

\begin{verbatim}
initialize x
initialize directions d_i

repeat:

    x_start = x

    for i in directions:

        alpha = line_minimize(f,x,d_i)

        x = x + alpha*d_i

    d_new = x - x_start

    replace direction with largest decrease

until convergence
\end{verbatim}

\end{frame}

%------------------------------------------------
\section{Python Example}

\begin{frame}[fragile]{Example Implementation}

\begin{verbatim}
import numpy as np
from scipy.optimize import minimize

def f(x):
    return x[0]**2 + 5*x[1]**2

x0 = np.array([3,4])

res = minimize(f, x0, method='Powell')

print(res.x)
\end{verbatim}

Uses SciPy's Powell implementation.
\end{frame}

%------------------------------------------------
\section{Applications}

\begin{frame}{Applications in Physics}

Powell's method used in:

\begin{itemize}
\item Variational parameter optimization
\item Molecular geometry optimization
\item Monte Carlo parameter fitting
\end{itemize}
\end{frame}

%------------------------------------------------
\begin{frame}{Example}

Minimization of

\[
f(x,y) = x^2 + 5y^2
\]

Powell method converges in few iterations

because function is quadratic.
\end{frame}

%------------------------------------------------
\section{Summary}

\begin{frame}{Summary}

\begin{itemize}
\item Powell's method is a derivative-free optimization algorithm
\item Uses sequential line minimizations
\item Updates search directions dynamically
\item Approximates conjugate directions
\item Efficient for moderate dimensional smooth problems
\end{itemize}
\end{frame}

%------------------------------------------------
\begin{frame}{References}

\begin{itemize}

\item M. J. D. Powell (1964)

\item Numerical Recipes

\item Nocedal \& Wright
Numerical Optimization

\end{itemize}

\end{frame}

\end{document}


⸻

What this lecture contains

The slides include

Theory
	•	motivation for derivative-free optimization
	•	geometry of line search methods
	•	conjugate directions concept
	•	Powell derivation for quadratic functions

Algorithm
	•	step-by-step Powell iteration
	•	direction updates
	•	convergence conditions

Implementation
	•	pseudocode
	•	Python example

Applications
	•	variational physics optimization
	•	parameter fitting

⸻

If you want, I can also produce a much more advanced lecture (80–100 slides) including:
	•	full geometric derivation of conjugate directions
	•	comparison with conjugate gradient and BFGS
	•	Brent line search derivation
	•	numerical examples
	•	visualization of optimization paths
	•	Python code implementing Powell from scratch

That version is often used in graduate computational physics or machine learning courses.
